
\section{A lower bound at fixed volume, and why it does not reach Hypothesis (A)}\label{sec:smallbox}

Every result proved in this paper bounds the gap from above. This section records one lower bound
at weak coupling. It holds at fixed spatial volume, namely one spatial site, and in one symmetry sector.
Its proof was written by a session of a second model lineage (OpenAI Codex), not by the sessions that
wrote this paper, and has not been read by a person. We report it with its provenance and explain why it gives nothing for
Hypothesis~\ref{hyp:A}.

\subsection{Setting}
Fix one spatial site with periodic identification in $k=2$ or $k=3$ spatial directions, and infinite
time. The spatial links are $U_1,\dots,U_k\in\SU(2)$. Write a unit quaternion as $U=(u_0,\mathbf u)$ with
$\mathbf u\in\R^3$ its imaginary part. The spatial plaquettes give the potential
\[
V=2\beta\sum_{i<j}|\mathbf u_i\times\mathbf u_j|^2 .
\]
Let $K$ be the product of independent convolutions, one per link, by the Wilson increment whose
density with respect to probability Haar measure is proportional to $e^{\beta\operatorname{Re}\Tr U/2}$,
normalised so that $K1=1$. The transfer operator is $T=e^{-V/2}Ke^{-V/2}$ on $L^2(\SU(2)^k)$. With the
unnormalised kernel $e^{\beta(\frac12\operatorname{Re}\Tr U-1)}$, $-\log T$ is shifted by a
$\beta$-dependent constant; gaps are unaffected, absolute energies are not.

The \emph{center-even sector} consists of the functions invariant under simultaneous conjugation of all
links (gauge invariance at the single site) and even under each independent center flip
$U_i\mapsto-U_i$. On it the temporal gauge projector acts as the identity. Let
$\lambda_0\ge\lambda_1\ge\cdots$ be the eigenvalues of $T$ there, counted with multiplicity, and
$h=\beta^{-1/3}$. The limiting operator is
\[
h_k=-\tfrac12\sum_{i=1}^k\Delta_i+2\sum_{i<j}|x_i\times x_j|^2
\qquad\text{on the rotation-invariant subspace of }L^2((\R^3)^k),
\]
rotations acting simultaneously on $x_1,\dots,x_k$. Its eigenvalues, ordered and counted with
multiplicity, are $e_0\le e_1\le\cdots$. The potential vanishes on the commuting configurations
(all $x_i$ parallel), the valley.

\subsection{The reported statement}

\medskip\noindent\textbf{Reported} (one-site limit in the center-even sector). \emph{For $k=2$ or $3$ and
each fixed $j$, as $\beta\to\infty$,}
\[
\frac{1-\lambda_j}{h}\longrightarrow e_j,\qquad
\frac{\log(\lambda_0/\lambda_1)}{h}\longrightarrow e_1-e_0>0,\qquad h=\beta^{-1/3}.
\]

\noindent\emph{Status.} Written proof by a session of a second model lineage (OpenAI Codex);
checked step by step by a session of the first lineage (Claude), which found $27$ of $30$ items
standing, none defective and $3$ not checked, and one wording slip in the companion proof of the vacuum
energy; not read by a person; the Lean files cover scalar lemmas only. It is not a theorem of this paper and is
not proved here. The full text is in the research repository named in~\cite{Repo} (a different folder from
the one cited there), folder \url{yang-mills/observable-normalization/one-site-vacuum/} and its
subfolders, at commits \repo{c5c1ea9} (vacuum energy), \repo{4966c89} (the statement above) and
\repo{dd2927a} (Remark~\ref{sb:rem-center}).
The proof does not give a numerical value for $e_1-e_0$. A companion statement by the same author
gives $c\beta^{-1/3}\le-\log\lambda_0\le C\beta^{-1/3}$ for the vacuum energy alone.

\emph{Proof sketch} (as written by its author; numbering of the steps ours).
\begin{enumerate}
\item \emph{Fibre inequality.} Fix $i$, pick $j\ne i$, freeze every link except $U_j$, and let
$r=|\mathbf u_i|\ge h$. Then $I-K_j+V_{ij}\ge c\,r$ uniformly in the frozen variables. The proof runs the
one-link walk for $m=\lceil 8/r\rceil$ steps, uses the smoothing estimate of the vacuum-energy proof and a
tube estimate for the transverse component, and bounds the survival of the walk killed at rate $V_{ij}$.
Since $K\le K_j$ and $V\ge V_{ij}$, this gives
$c\int|\mathbf u_i|\,|g|^2\le\langle g,(I-K+V)g\rangle+ch\|g\|^2$ for every $g$.
\item \emph{Compactness.} For a unit vector $F$ with $\langle F,(I-T)F\rangle\le Ch$ put $g=e^{-V/2}F$.
Step~1 bounds the mass of $g$ outside $\{\max_i\theta_i\le Rh\}$ ($\theta_i$ the distance of $U_i$ to
the nearer center) by $C'/R$. This is where the valley arms are handled: Step~1 penalises every link
away from a center whatever the other links do, including configurations with a common axis, so the
escape along the commuting valley is controlled. A kinetic bound $1-K\ge c'\min(D/\beta,1)$ by Bessel
ratios ($D$ the product Laplacian), an isometry $J_h$ onto Euclidean $L^2$ in the scaled variables
$U_i=\exp(hx_i)$, and Rellich's theorem on balls give strong $L^2$ compactness of $J_hF$.
\item \emph{Consistency on test functions.} For smooth compactly supported invariant $\psi$,
$J_h(I-T)J_h^*\psi/h\to h_k\psi$ in $L^2$: the potential from $V/h$, the Laplacian from the second
moment of one increment ($\beta\,\E\,\xi_a\xi_b\to\delta_{ab}$), with a third-order remainder
$O(\beta^{-3/2}h^{-3})=O(h^{3/2})$.
\item \emph{Min--max.} Trial spaces from Step~3 give $\limsup(1-\lambda_j)/h\le e_j$. Eigenvectors of
$(I-T)/h$ have strong limits by Step~2 that remain orthonormal and are weak eigenfunctions of $h_k$, which
gives the $\liminf$. The operator $h_k$ has compact resolvent and a simple positive ground state, so
$e_1>e_0$; finally $-\log\lambda_j=(1-\lambda_j)+O(h^2)$.
\end{enumerate}

A second proof of the case $k=2$, with a different architecture (a partition into four center boxes and
capped valley arms in place of the walk estimate), was written by sessions of the first lineage and
Its upper half and its lemmas up to tightness were challenged once by a session of the same lineage
($11$ of $12$ items stood; one sign defect, the bound used survives). The concluding steps of the lower half,
written in full by another session of that lineage, which also rederived one constant and replaced one
partition of unity, were challenged once in the same way ($22$ of $24$ items stood, one not checked; two
incomplete justifications and one overestimated exponent, none affecting the conclusion). It has not been read by a person, and it is not reproduced here.

\begin{remark}[All center sectors]\label{sb:rem-center}
With all $2^k$ center characters restored to the gauge-invariant space, the same author reports that for
some $0<c<C$ depending on $k$, for all large $\beta$ and every nontrivial character $\chi$,
$e^{-C\sqrt\beta}\le\log(\lambda_0/\lambda_\chi)\le e^{-c\sqrt\beta}$, where $\lambda_\chi$ is the largest
eigenvalue in the sector $\chi$. The first gap on the full gauge-invariant space is then between
$e^{-C\sqrt\beta}$ and $e^{-c\sqrt\beta}$ for large $\beta$, and the $\beta^{-1/3}$ statement above holds in
the center-even sector only. Status: written proof by a session of the second lineage; it was not part of
the step-by-step check above; not read by a person. Measured center-odd splittings for $k=2$ (double precision up to $\beta=64$, double-double arithmetic up to
$\beta=256$) are fitted by $-\log m\approx c\sqrt\beta+q\beta^{1/6}$ with $c\approx4.60$; a heuristic gives
$c=\int_0^\pi\sqrt{2\operatorname{arccosh}(1+2\sin^2a)}\,da=4.6054$. Six points do not establish the form. Sessions of the first lineage have written a proof
of one side, $\liminf_{\beta\to\infty}(-\log m)/\sqrt\beta\ge4.6054$ for the singly odd sector at $k=2$ (challenged
within that lineage, not read by another lineage or by a person, not reproduced here); a proof of the
opposite inequality has since been assembled from separately written parts and challenged within the same
lineage; it is held to the same standing and is not relied on here.
\end{remark}

\subsection{Numbers}
Table~\ref{sb:tab-num} collects exact numerics for both models. All entries are measured, none is proved.
The lattice values are logarithms of ratios of transfer eigenvalues in the center-even sector, in lattice
units; the continuum values are Rayleigh--Ritz computations for $h_k$.

\begin{table}[ht]
\centering
\small
\begin{tabular}{@{}p{0.43\textwidth}p{0.49\textwidth}@{}}
\toprule
Quantity & Value (status: measured) \\
\midrule
\multicolumn{2}{@{}l}{\emph{Two directions}} \\
$e_0$, $e_1$ of $h_2$ & $3.3505330$, $5.3254970$; $e_1-e_0=1.974964$ \\
first scalar excited level of $h_2$ & $5.725408$ (gap $2.374875$) \\
$\log(\lambda_0/\lambda_1)\,\beta^{1/3}$, $\beta=32,64,128$ & $1.636$, $1.760$, $1.837$ \\
$\log(\lambda_0/\lambda_1)\,\beta^{1/3}$, $\beta=256,512,1024$ & $1.887$, $1.919$, $1.939$ \\
two-term fit in $\beta^{-2/3}$, $\beta^{-1}$ & $1.97496$ (endpoint deletion: $1.97503$, $1.97489$) \\
\midrule
\multicolumn{2}{@{}l}{\emph{Three directions}} \\
$e_0$ of $h_3$ & $6.534885$ \\
first gap of $h_3$ (five-fold, $J=2$) & $3.012539$ \\
first scalar gap of $h_3$ & $3.602838$ \\
$e_0$ in the normalisation of~\cite{vB00} & $4.1167196$ (published: $4.116719735$) \\
lattice $\log(\lambda_0/\lambda_1)$, $\beta=1,2,4$ & $2.876117$, $1.621505$, $0.968930$ \\
the same times $\beta^{1/3}$ & $2.88$, $2.04$, $1.54$ (not asymptotic) \\
\bottomrule
\end{tabular}
\caption{Exact numerics for the one-site model, center-even sector. Measured, not proved.}\label{sb:tab-num}
\end{table}

\emph{Two directions.} The model has an $\mathrm{O}(2)$ symmetry mixing $x_1$ and $x_2$, and $e_1=e_2$ is an
exact doublet of charge two. That the lowest non-scalar branch is this doublet is proved; that it lies
below the first excited scalar level rests on the numerics. A third computation (a Legendre and oscillator basis)
reproduces $e_0$ and $e_1$ to eight digits and puts the first non-degenerate excited level at $5.725408$, $0.40$
above the doublet. By the heuristic of Section~\ref{sb:sub-why} with $L^3$ replaced by $L^2$, $h_2$ is also the
leading constant-mode Hamiltonian of the theory in $2+1$ dimensions on a small spatial torus; there the
$\mathrm{O}(2)$ is the rotation group of space, the doublet has spin two and the level at $5.725408$ is the scalar. The lattice symmetry splits the doublet
($0.19238$ against $0.19331$ at $\beta=1024$); the square group permits but does not force this splitting.
The lattice values approach $e_1-e_0$ from below; at $\beta=1024$, $(1-\lambda_j)/h$ is still $16$ to
$23$ per cent below $e_j$, while $-\log\lambda_j/h$ is within $1.3$ per cent.
The fitted limit $1.97496$ is a four-point fit with one residual degree of freedom; deleting either endpoint
moves it to $1.97503$ or $1.97489$. It is a comparison with $e_1-e_0$, not a certified limit. A three-point
geometric extrapolation gives $1.97583$.
The lowest center-even gap was computed by two methods (a Peter--Weyl diagonalization and a Nystr\"om
discretization with Lanczos), which agree to $4.6\cdot10^{-9}$ or better at $\beta=2,\dots,1024$; the second
level agrees to $1.1\cdot10^{-8}$.

\emph{Three directions.} The lattice model was converged only up to $\beta=4$; at $\beta=8$ it is not
converged, and the computed range is not asymptotic. The limit was therefore tested by Monte Carlo instead (32 independent chains per $\beta$, a variational analysis of center-even operators in the channels $E$, $T_2$, $A_1$ of the cubic group, plateau rule fixed in advance for the main runs; controls: the exact gaps at $\beta=2$ and $4$, which failed the rule fixed in advance and were reproduced only by analyses chosen afterwards (an earlier reference time; a differenced correlator removing a constant floor), and the two-direction value at $\beta=256$, which came out $1.8$ standard errors high). The products $\log(\lambda_0/\lambda)\,\beta^{1/3}$ are, for $E$, $T_2$, $A_1$: $2.747(10)$, $2.813(13)$, $3.225(36)$ at $\beta=64$; $2.955(9)$, $2.973(8)$, $3.524(17)$ at $\beta=1024$; $2.992(13)$, $2.989(13)$, $3.544(27)$ at $\beta=4096$; $2.988(23)$, $2.984(23)$, $3.568(48)$ at $\beta=16384$; against $3.0125$, $3.0125$, $3.6028$. The approach is from below, as for two directions, and the $E$--$T_2$ splitting falls to $0.018(11)$ at $\beta=1024$ and vanishes within errors at $\beta=4096$ and $16384$. This is a measurement at one to two per cent.
In the normalisation of the small-volume literature, $-\tfrac12\Delta+\tfrac12\sum_{i<j}|x_i\times x_j|^2$
(energies divided by $4^{1/3}$), $e_0=4.1167196$. The published value is $4.116719735$
\cite[text after eq.~(12)]{vB00}; the difference, $1.5\cdot10^{-7}$, is within the rounding of
$6.534885$. That source calls the number the lowest-order mass of the scalar glueball; in its own
eq.~(12), and here, it is a ground energy, not a gap. The two gaps, $1.89778$ and $2.26965$ in that
normalisation, compare with $1.897780$ and $2.269638$ obtained by subtracting the vacuum from the lowest
even tensor and the first excited even scalar levels of Pavel~\cite[Tables~1 and~2]{Pa07} (found by a session
of the second lineage and read by a session of the first); the differences are $3\cdot10^{-7}$ and
$7\cdot10^{-6}$. Our scalar gap is the less converged of the two, so $3.602838$ in Table~\ref{sb:tab-num} should
be read as $3.60283$. The original source~\cite{LM84} was reached as an abstract only; its leading scalar
coefficient is quoted as $2.2696390$ in a later paper that we did not read ourselves.

\subsection{Why it does not reach Hypothesis~\ref{hyp:A}}\label{sb:sub-why}
The following derivation is heuristic. The reported statement is for one spatial site ($L=1$) only.

On the periodic box of side $L$ restrict to constant spatial links $U_i=\exp(\tfrac{i}{2}A_i\cdot\sigma)$
and to small $A_i$. Each spatial plaquette gives $1-\tfrac12\Tr U_P=\tfrac18|A_i\times A_j|^2(1+O(A^2))$ and
there are $L^3$ of them per plane; the temporal plaquettes give the kinetic term
$\tfrac{\beta L^3}{8}\sum_i|\dot A_i|^2$. Treating the transfer operator as $e^{-H}$,
\[
H=-\frac{2}{\beta L^3}\Delta_A+\frac{\beta L^3}{8}\sum_{i<j}|A_i\times A_j|^2,
\qquad A=2(\beta L^3)^{-1/3}Y \ \Longrightarrow\ H=(\beta L^3)^{-1/3}h_3 .
\]
So the constant modes predict a gap $m\approx e\,\beta^{-1/3}/L$ in lattice units, with $e=3.012539$ the
first gap of $h_3$ (a measured number). At fixed $\beta$ this tends to zero as $L\to\infty$.
With $\beta=4/g^2$ this is $mL=e'g^{2/3}$, $e'=4^{-1/3}e$: the $g^{2/3}/L$ form of the small-volume
expansion of L\"uscher~\cite{Lu83}, with the leading-order effective Hamiltonian of~\cite[eq.~(9)]{vB00}
(the agreement of coefficients was checked against~\cite{vB00}, not against~\cite{Lu83}); the bare coupling $g$ here agrees with the
renormalised $g(L)$ of the literature at leading order only. The derivation neglects the non-constant
modes and assumes the replacement of $T$ by $e^{-H}$ on the box, the box analogue of the reported
statement. For each fixed box the same author has since written a derivation of this scaling, with
coefficient $(e_1-e_0)/L$ on a cube of side $L$ and constants not uniform in $L$ (repository folder
\url{yang-mills/observable-normalization/periodic-box-spectrum/}, commit \repo{d0a7ccb}). A session of the
first lineage read it and found no false line; it found the two-sided statement for the full spectrum
conditional on an exterior estimate it could not check, and two gaps of exposition. An exact Gaussian
computation on boxes $2^3$, $3\times2\times2$ and $3^3$ confirmed the factor $L^3$ in the mass of the constant
mode. A Monte Carlo measurement on $2^3\times T$ ($16$ chains, differenced correlators, analysis fixed in advance) gives for
$\log(\lambda_0/\lambda)\,\beta^{1/3}$ in the channels $E$, $T_2$: $1.497(19)$, $1.507(16)$ at $\beta=256$ and $1.502(24)$,
$1.509(18)$ at $\beta=1024$, against the predicted $3.012539/2=1.5063$; the scalar channel ($1.900(24)$, $1.890(36)$
against $1.8014$) is not resolved, its effective mass still falling after the accepted window. We do not
rely on that derivation: the conclusion of this subsection needs only that the predicted gap vanishes as
$L\to\infty$.
The literature on this expansion places the end of its regime at $z=mL$ of order one: effects of the
boundary of the fundamental domain from $z>0.9$ \cite[pp.~11--12]{vB00}, and a crossover near $z\approx2$
(\cite{LM84}, abstract). The route therefore gives no lower bound uniform in the volume, and no lower bound
of the kind Hypothesis~\ref{hyp:A} requires. It is the lattice form of a known regime, at one site.

In a limited search we did not find a proof, in the literature, of the discrete-time one-site limit stated
above, nor a theorem giving a lower bound on the gap on a finite torus at weak coupling that is uniform in
any parameter, nor a published spectrum of $h_2$. The searches were limited (catalogue phrase queries
returned nothing even where they should have), so ``not found'' is weak evidence.
Of the sources, \cite{vB00} was read in full; \cite{Lu83} and~\cite{LM84} were read as abstracts only.

\subsection{A false start}
For $k=2$ the first form of coercivity suggested in this work was
$1-K+V\ge|\sin a_1|+|\sin a_2|-\beta^{-1/3}$ ($a_i$ the class angles), with coefficient one on the weight.
It is false: $1-K\le1$, while $V$ is small on high-spin states concentrated near $a_1=a_2=\pi/2$ with a
common axis, where the weight equals $2$; numerically the infimum of the spectrum of the difference, times
$\beta^{1/3}$, was below $-2.39$ at $\beta=16$ and below $-3.35$ at $\beta=64$. A weight of this kind must be
capped, or carry a small coefficient as in Step~1 of the sketch above; the second proof uses a capped fibre
lemma.

